% =====================================================================
% 別冊『Git & GitHub 実践 — チームで安全に開発する』
% エンジン: LuaLaTeX（jlreq / book）
% ビルド:   latexmk main_team.tex   （出力は output/main_team.pdf。エンジン指定は .latexmkrc 側）
%           全冊まとめて作るときは、このフォルダで latexmk のみ実行
%
% 対象: 本編『はじめてのバージョン管理』を終え、チーム開発に進む人。
% ブランチ・プルリクエスト（PR）・コンフリクト解決を扱う。
% =====================================================================
\documentclass[book, paper=a4, fontsize=10.5pt, titlepage]{jlreq}

% =====================================================================
% preamble.tex — パッケージ読込・色・カスタム環境・マクロの定義
% 本文（chapters/*.tex）から使う見た目の部品はすべてここに集約する。
% エンジン: LuaLaTeX + jlreq
% =====================================================================

% --- 基本パッケージ ---------------------------------------------------
\usepackage{graphicx}          % 画像の取り込み
\usepackage{xcolor}            % 色指定
\usepackage{booktabs}          % 罫線のきれいな表
\renewcommand{\arraystretch}{1.25} % 表の行を少しゆったりさせる
\usepackage{tabularx}          % 幅を自動調整する表
\usepackage{enumitem}          % 箇条書きの調整
\usepackage{fancyvrb}          % コード枠に使う Verbatim
\usepackage[most]{tcolorbox}   % 装飾枠（コード・注意書き）
\usepackage{adjustbox}         % 図が本文幅を超えたときの自動縮小
\usepackage{tikz}              % 概念図・フロー図
\usetikzlibrary{arrows.meta,positioning,shapes.geometric,calc,fit,backgrounds}
\usepackage{float}             % 図を本文の位置に固定する [H]
\usepackage{caption}           % 図番号つきキャプションの体裁
\captionsetup{font=small, labelfont=bf, labelsep=quad}

% 画像の探索パス（既存の Markdown 版と同じ images/ を共有する）
\graphicspath{{../images/}{./images/}}

% --- 配色 -------------------------------------------------------------
% Markdown 版の Mermaid 図の色をそのまま踏襲する（薄い色は使わない方針）。
\definecolor{gitRed}{HTML}{C0392B}     % 作業ツリー・注意
\definecolor{gitOrange}{HTML}{D68910}  % ステージング
\definecolor{gitGreen}{HTML}{2E7D32}   % リポジトリ・完了
\definecolor{gitBlue}{HTML}{4A90D9}    % コミット・情報
\definecolor{gitPurple}{HTML}{6C3483}  % 自分側
\definecolor{gitTeal}{HTML}{117A65}    % 両方採用
\definecolor{gitAmber}{HTML}{8A6D0B}   % 判断ノード
\definecolor{codeBg}{HTML}{F4F5F7}     % コード枠の背景
\definecolor{codeFrame}{HTML}{D0D3D9}  % コード枠の枠線

% --- コード枠 ---------------------------------------------------------
% 使い方: \begin{codebox}[bash] ... Verbatim ... \end{codebox}
% 省略時のラベルは bash。右上に言語ラベルを小さく表示する。
\newtcolorbox{codebox}[1][bash]{%
  enhanced, breakable,
  colback=codeBg, colframe=codeFrame, boxrule=0.6pt, arc=2pt,
  left=9pt, right=9pt, top=13pt, bottom=8pt,
  before skip=8pt, after skip=8pt,
  overlay={%
    \node[anchor=north east, font=\sffamily\scriptsize\bfseries,
          text=gray!60] at ([xshift=-6pt, yshift=-3pt]frame.north east) {#1};%
  },
}
% コード本体を書くための Verbatim 環境（小さめ・行間やや詰め）
\DefineVerbatimEnvironment{code}{Verbatim}{fontsize=\small, xleftmargin=0pt}

% 本文中のインラインコード: \code{git status}
\newcommand{\cmd}[1]{\texttt{#1}}

% --- 注意書き（admonition）枠 ----------------------------------------
% ヒント（緑）・注意（赤）・補足（グレー）の3種。emoji は使わずラベルで表す。
% タイトルは本文（明るい背景）に直接乗るため、白文字だと見えない。
% タイトル文字色は各テーマ色にして、はっきり読めるようにする。
\newtcolorbox{hintbox}{%
  enhanced, breakable,
  colback=gitGreen!5, colframe=gitGreen, boxrule=0.8pt, arc=2pt,
  left=8pt, right=8pt, top=6pt, bottom=6pt,
  coltitle=gitGreen, fonttitle=\sffamily\bfseries\small,
  title=ヒント, attach title to upper=\quad,
}
\newtcolorbox{warnbox}{%
  enhanced, breakable,
  colback=gitRed!5, colframe=gitRed, boxrule=0.8pt, arc=2pt,
  left=8pt, right=8pt, top=6pt, bottom=6pt,
  coltitle=gitRed, fonttitle=\sffamily\bfseries\small,
  title=注意, attach title to upper=\quad,
}
% 執筆メモ（雛形用の枠）。原稿を清書したら、この枠ごと削除する。
\newtcolorbox{draftbox}{%
  enhanced, breakable,
  colback=gitAmber!6, colframe=white, boxrule=0pt, arc=2pt,
  borderline={0.9pt}{0pt}{gitAmber, dashed},
  left=8pt, right=8pt, top=6pt, bottom=6pt,
  coltitle=gitAmber, fonttitle=\sffamily\bfseries\small,
  title=執筆メモ, attach title to upper=\quad,
}
\newtcolorbox{notebox}{%
  enhanced, breakable,
  colback=gray!8, colframe=gray!60!black, boxrule=0.8pt, arc=2pt,
  left=8pt, right=8pt, top=6pt, bottom=6pt,
  coltitle=gray!30!black, fonttitle=\sffamily\bfseries\small,
  title=補足, attach title to upper=\quad,
}

% --- スクリーンショット取り込み --------------------------------------
% \screenshot[幅の割合]{ファイル}{キャプション}
% フロートにせず本文の流れに沿って中央寄せで置く（説明の直後に見せたいため）。
\newcommand{\screenshot}[3][0.7]{%
  \begin{center}
    \includegraphics[width=#1\linewidth]{#2}\\[3pt]
    {\small\itshape #3}%
  \end{center}%
}

% \screenshottrim[幅の割合]{ファイル}{キャプション}{trim指定}
% 余白の多いスクリーンショットを切り詰めて大きく見せる。
% trim は「左 下 右 上」の順で、画像1ピクセル = 1bp として指定する
% （例: 1920x1080 の下半分を捨てる → {0 540bp 0 0}）。
\newcommand{\screenshottrim}[4][0.95]{%
  \begin{center}
    \includegraphics[width=#1\linewidth, trim=#4, clip]{#2}\\[3pt]
    {\small\itshape #3}%
  \end{center}%
}

% --- TikZ 共通スタイル ------------------------------------------------
% Mermaid のフロー図を再現するためのノード／矢印スタイル。
\tikzset{
  flownode/.style={%
    rounded corners=4pt, draw=none, text=white,
    font=\sffamily\small, align=center,
    minimum height=1.15cm, minimum width=2.2cm, inner sep=6pt},
  decisionnode/.style={%
    diamond, aspect=1.9, draw=none, text=white,
    font=\sffamily\small, align=center, inner sep=3pt},
  flowarrow/.style={-{Latex[length=3.6mm, width=3mm]}, line width=1.8pt, draw=black},
  edgelabel/.style={font=\sffamily\normalsize\bfseries, fill=white, inner sep=3pt},
}
% 図全体を中央寄せで置くためのヘルパ環境（図番号なし）。
% 本文幅を超えた図は自動で縮小する（幅が足りている図は拡大しない）。
\newenvironment{diagram}%
  {\par\medskip\begin{center}\begin{adjustbox}{max width=\linewidth}}%
  {\end{adjustbox}\end{center}\par\medskip}

% --- 図番号つきの図（本文が「図1.2のように」と参照する概念図・フロー図用）---
% 手順に紐づく操作スクショには使わない（そちらは \screenshot を使う）。
\makeatletter
% TikZ など本体を囲む環境: \begin{numfig}{キャプション}{ラベル} ... \end{numfig}
\newenvironment{numfig}[2]{%
  \def\nf@cap{#1}\def\nf@lab{#2}%
  \begin{figure}[H]\centering\begin{adjustbox}{max width=\linewidth}%
}{%
  \end{adjustbox}\caption{\nf@cap}\label{\nf@lab}\end{figure}%
}
\makeatother
% 画像1枚を図番号つきで置く: \figimage[幅]{ファイル}{キャプション}{ラベル}
\newcommand{\figimage}[4][0.95]{%
  \begin{figure}[H]\centering
    \includegraphics[width=#1\linewidth]{#2}
    \caption{#3}\label{#4}
  \end{figure}%
}


% hyperref はプリアンブルの最後（他パッケージ読込後）に置くのが安全。
\usepackage[unicode, hidelinks]{hyperref}

% --- 書誌情報 ---------------------------------------------------------
\title{Git \& GitHub 実践}
\author{}
\date{\today}

\begin{document}

% --- 表紙 -------------------------------------------------------------
\begin{titlepage}
  \centering
  \vspace*{4cm}
  {\sffamily\bfseries\Huge Git \& GitHub 実践\par}
  \vspace{1.2cm}
  {\sffamily\Large チームで安全に開発する\par}
  \vspace{0.6cm}
  {\sffamily\normalsize ── 本編『はじめてのバージョン管理』の続き ──\par}
  \vfill
  {\large \today\par}
\end{titlepage}

% --- 目次 -------------------------------------------------------------
\tableofcontents
\clearpage

% --- 本文 -------------------------------------------------------------
% =====================================================================
% 別冊 第1章 ブランチとPR編 — チームで安全に開発する
% 元原稿: 03_ブランチ・PR編.md（前半：ブランチ〜PR〜実践フロー）
% =====================================================================
\chapter{ブランチとPR編 —チームで安全に開発する—}

本書は本編『はじめてのバージョン管理』の続きです。本編では、ローカルとGitHubをつなぎ、\textbf{1本の流れ（\cmd{main}）}をアップロード／ダウンロードしました。本書では、\textbf{ブランチ}で作業を枝分かれさせ、\textbf{プルリクエスト（PR）}でレビューを挟み、\textbf{コンフリクト（競合）}を解決します。つまり「複数人・複数タスクを安全に並行させる」方法を扱います。

\begin{hintbox}
\textbf{この本の読み方}
\begin{itemize}[leftmargin=1.4em, nosep]
  \item 1人で使う分には本編で十分です。チームで開発する、または作業を安全に分けたくなったらここを読みます
  \item チームの運用ルールやPRの書き方に迷ったら → 付録A「おすすめのブランチ運用ルールとPRの書き方」へ
  \item コマンドをド忘れした人 → 『逆引きチートシート』へ
  \item エラーが出て困った人 → 『トラブルシューティング』へ
\end{itemize}
\end{hintbox}

% ---------------------------------------------------------------------
\section{なぜブランチが必要か}

これまでは \cmd{main} という1本の流れだけで作業してきました。\cmd{main} に直接編集を重ねると、こんな問題が起きます。

\begin{itemize}[leftmargin=1.4em]
  \item \textbf{未完成の変更が本流に混ざる}：書きかけのコードや原稿が \cmd{main} を壊す
  \item \textbf{複数人が同じ流れを触って衝突する}：2人が同時に \cmd{main} を編集すると \cmd{push} がぶつかる
  \item \textbf{お試しがやりにくい}：「ダメなら元に戻す」を安全に試せない
\end{itemize}

\textbf{ブランチ}を使うと、\cmd{main} から枝分かれした作業用の流れを作れます。作業はその枝の上で行い、完成してから \cmd{main} に合流（マージ）させます。こうすれば \cmd{main} は常に「動く・完成した状態」に保たれます（図\ref{fig:branch-merge}）。

\begin{numfig}{ブランチで枝分かれし、完成してから main に合流させる}{fig:branch-merge}
\begin{tikzpicture}
  \node[flownode, fill=gitGreen] (M0) {main};
  \node[flownode, fill=gitGreen, right=2.4cm of M0] (M1) {main};
  \node[flownode, fill=gitGreen, right=4.8cm of M1] (M2) {main\\（マージ後）};
  \node[flownode, fill=gitOrange, above=1.4cm of M1] (F1) {feature\\作業中};
  \node[flownode, fill=gitOrange, right=2.4cm of F1] (F2) {feature\\完成};
  \draw[flowarrow] (M0) -- (M1);
  \draw[flowarrow] (M1) -- (M2);
  \draw[flowarrow] (F1) -- (F2);
  \draw[flowarrow, dashed] (M0) to[out=70, in=200]
        node[edgelabel, midway, left] {枝分かれ} (F1);
  \draw[flowarrow, dashed] (F2) to[out=-20, in=110]
        node[edgelabel, midway, right] {マージ} (M2);
\end{tikzpicture}
\end{numfig}

% ---------------------------------------------------------------------
\section{ブランチとは}

\textbf{ブランチ（branch＝枝）}は、その名の通り履歴の枝分かれです。ある時点の \cmd{main} からコピーを取り、そこから独立して作業を進める「もう1本の作業ライン」です。

\begin{itemize}[leftmargin=1.4em]
  \item \cmd{main}（メインブランチ）… 完成した状態を保つ本流。通常ここは直接編集しない
  \item \cmd{feature/...} などの作業ブランチ … 1つのタスク専用の枝。ここで自由に編集・コミットする
\end{itemize}

各ブランチは独立しているので、作業ブランチで何をしても \cmd{main} には影響しません。\cmd{git switch} でブランチを切り替えると、作業ツリーの中身もそのブランチの内容に入れ替わります。

分け方の基本は\textbf{「1タスク＝1ブランチ」}です。「ログイン機能の追加」「READMEの修正」のように、1つの目的ごとにブランチを分けると、後のPRやレビューがきれいにまとまります。

\subsection{たとえ：倉庫で考えるブランチ}

本編では、リポジトリを書類を保管する「倉庫」にたとえました。ブランチは、この倉庫の中に置かれた\textbf{棚}だと考えると分かりやすくなります。

\begin{center}
\begin{tabularx}{\linewidth}{@{}l X@{}}
\toprule
Git & 倉庫のたとえ \\
\midrule
\cmd{main} & 全員が共有する正式な棚（本棚）。常に「完成品」だけを並べる \\
作業ブランチ & 本棚の中身を丸ごと写し取った、自分専用の作業棚 \\
そのブランチでの \cmd{commit} & 自分の作業棚に、書き直した書類の写しを並べていく \\
PR → マージ & 完成した書類を本棚に納める申請と反映 \\
\bottomrule
\end{tabularx}
\end{center}

作業ブランチを作ると、\cmd{main} の棚に並んでいる書類束（コミット）をそのまま写し取った、\textbf{自分専用の棚がもう1つ}用意されます。倉庫（リポジトリ）は1つのままで、\textbf{その中に棚が2つ並ぶ}というのが、ブランチを作るということです（図\ref{fig:branch-create}）。

\figimage[0.95]{03_branch・PR/create_branch.png}{ブランチを作ると、倉庫の中に \texttt{main} の棚と作業ブランチの棚が並ぶ}{fig:branch-create}

作業棚は本棚の写しから始まる独立した棚なので、そこで何をしても本棚（\cmd{main}）は汚れません。だから安心して書きかけの書類を置いたり、失敗した写しを捨てたりできます。

\subsection{複数人のときはどうなるか}

このたとえは、複数人での作業でこそ効いてきます。メンバーはそれぞれ、本棚から自分専用の作業棚を写し取って並行して作業します。互いの作業棚は独立しているので、他の人の手元を邪魔しません（図\ref{fig:multi-branch}）。

\begin{numfig}{各メンバーが本棚から作業棚を写し取り、完成したら本棚へ納める}{fig:multi-branch}
\begin{tikzpicture}
  \node[flownode, fill=gitGreen, text width=2.6cm] (MAIN)
        {本棚（main）\\完成品の共有棚};
  \node[flownode, fill=gitOrange, text width=3cm, above right=1.9cm and 5cm of MAIN] (A)
        {Aさんの作業棚\\feature/login};
  \node[flownode, fill=gitPurple, text width=3cm, below right=1.9cm and 5cm of MAIN] (B)
        {Bさんの作業棚\\feature/report};
  \draw[flowarrow, dashed] (MAIN.north east) to[out=55, in=210]
        node[edgelabel, pos=0.62, above left=-1pt] {写し取る} (A.west);
  \draw[flowarrow] (A.south west) to[out=250, in=30]
        node[edgelabel, pos=0.4, below right=-1pt] {PR→マージ} (MAIN.east);
  \draw[flowarrow, dashed] (MAIN.south east) to[out=-55, in=150]
        node[edgelabel, pos=0.62, below left=-1pt] {写し取る} (B.west);
  \draw[flowarrow] (B.north west) to[out=110, in=-30]
        node[edgelabel, pos=0.4, above right=-1pt] {PR→マージ} (MAIN.east);
\end{tikzpicture}
\end{numfig}

完成したら、各自が自分の作業棚の書類を順番に本棚へ納めて（マージして）いきます。ここで、AさんとBさんが「同じ書類の同じ箇所」を別々に書き換えていた場合、後から納める人のときに「どちらの内容を正とするか決められない」状態になります。これが次章の\textbf{コンフリクト（競合）}です。

逆に言えば、\textbf{触るファイルや箇所をタスクごとに分けておく}と、本棚に納める順番が前後してもぶつかりません。これがコンフリクト予防の基本です。

% ---------------------------------------------------------------------
\section{ブランチの基本操作}

作業ブランチを作って切り替えるところから始めます。

\begin{codebox}[bash]
\begin{code}
# 現在のブランチ一覧を見る（* が今いるブランチ）
git branch

# ブランチを作って、同時にそこへ切り替える
git switch -c feature/add-login

# （参考）古い書き方。同じ意味
git checkout -b feature/add-login
\end{code}
\end{codebox}

\cmd{git switch -c} の \cmd{-c}（create）は「新しく作って切り替える」という意味です。すでにあるブランチへ移動するだけなら \cmd{-c} は不要です。

\begin{codebox}[bash]
\begin{code}
# 既存のブランチへ切り替える
git switch main
git switch feature/add-login

# 不要になったブランチを削除する（マージ済みのものだけ安全に消える）
git branch -d feature/add-login
\end{code}
\end{codebox}

作業ブランチに切り替えたら、あとは本編の黄金フローと同じです。「編集 → \cmd{git add} → \cmd{git commit}」を、このブランチの上で繰り返します。

\begin{hintbox}
\textbf{ブランチ名の付け方}：厳密な決まりはありませんが、\cmd{種類/内容} の形が広く使われます。例：新機能なら \cmd{feature/add-login}、バグ修正なら \cmd{fix/typo-in-readme}。目的がひと目で分かる名前にします。おすすめの命名規則やチームの運用ルールは、付録A「おすすめのブランチ運用ルールとPRの書き方」にまとめています。
\end{hintbox}

\begin{warnbox}
作業を始める前に、\textbf{\cmd{main} を最新にしてから枝分かれ}するのが安全です（\cmd{git switch main} → \cmd{git pull} → \cmd{git switch -c 新ブランチ}）。古い \cmd{main} から分岐すると、後でコンフリクトが起きやすくなります。
\end{warnbox}

% ---------------------------------------------------------------------
\section{ブランチをリモートへ上げる}

作った作業ブランチは、ローカルにあるだけではチームに共有されません。本編の \cmd{push} でリモートにも同じブランチを作ります。

\begin{codebox}[bash]
\begin{code}
# 初回：リモートに同名のブランチを作り、宛先を記憶させる
git push -u origin feature/add-login

# 2回目以降はこれだけでOK
git push
\end{code}
\end{codebox}

\cmd{-u}（\cmd{-{}-set-upstream}）は本編と同じで、「このローカルブランチの写し先はリモートのこのブランチだ」と結び付けるオプションです。一度付ければ、以降は \cmd{git push} / \cmd{git pull} だけで今いるブランチが同期されます（図\ref{fig:push-branch}）。

\begin{numfig}{作業ブランチをリモートにも作って同期する}{fig:push-branch}
\begin{tikzpicture}
  \node[flownode, fill=gitOrange, text width=2.8cm] (LB)
        {ローカル\\feature/add-login};
  \node[flownode, fill=gitBlue, text width=2.8cm, right=8cm of LB] (RB)
        {GitHub\\origin/feature/add-login};
  \draw[flowarrow] (LB) -- (RB)
        node[edgelabel, midway, above] {git push -u origin feature/add-login};
\end{tikzpicture}
\end{numfig}

\begin{hintbox}
リモートに上げたブランチは、GitHubのリポジトリ画面のブランチ切り替えメニューから確認できます。これで、他のメンバーもあなたのブランチを見られます。
\end{hintbox}

% ---------------------------------------------------------------------
\section{プルリクエスト（PR）を作る}

\textbf{プルリクエスト（Pull Request、PR）}は、「自分の作業ブランチの変更を \cmd{main} に取り込んでください」という依頼です。ただの合流ではなく、\textbf{レビュー（変更内容の確認）を挟む入口}になります（図\ref{fig:pr-flow}）。

\begin{numfig}{PRは push からマージまでにレビューを挟む}{fig:pr-flow}
\begin{tikzpicture}
  \node[flownode, fill=gitOrange, text width=2.4cm] (B) {feature を push};
  \node[flownode, fill=gitPurple, right=2.2cm of B] (P) {PRを作成};
  \node[flownode, fill=gitBlue, right=2.2cm of P, text width=2.2cm] (R) {レビュー\\（差分を確認）};
  \node[flownode, fill=gitGreen, right=2.2cm of R] (M) {main へマージ};
  \draw[flowarrow] (B) -- (P);
  \draw[flowarrow] (P) -- (R);
  \draw[flowarrow] (R) -- (M);
\end{tikzpicture}
\end{numfig}

\noindent\textbf{PRの作成手順}

\begin{enumerate}[leftmargin=1.6em]
  \item ブランチを \cmd{push} した後、GitHubのリポジトリ画面を開きます。上部に \textbf{「Compare \& pull request」} ボタンが出るのでクリックします（出ない場合は \textbf{Pull requests} タブ → \textbf{New pull request}）。
  \item \textbf{base}（取り込み先＝通常 \cmd{main}）と \textbf{compare}（自分の作業ブランチ）が正しいか確認します。
  \item \textbf{タイトル}と\textbf{説明}を書きます。「何を・なぜ変えたか」を簡潔に書くとレビューが早く進みます。
  \item \textbf{Create pull request} をクリックします。
\end{enumerate}

\screenshot[0.62]{03_branch・PR/PR_Create_Window.png}{PR作成画面（Compare \& pull request）}
\screenshot[0.62]{03_branch・PR/Create_PR_Setting.png}{base と compare を確認し、タイトル・説明を記入して作成する}

\begin{hintbox}
\textbf{説明の書き方が伝わりやすさを決めます。} 差分を見れば「どう変えたか」は分かるので、説明には\textbf{「なぜ変えたか」}を書くのがコツです。そのまま使えるPRテンプレートと記入例を、付録A「おすすめのブランチ運用ルールとPRの書き方」に用意しています。
\end{hintbox}

PRを作ると、\textbf{以降その作業ブランチに \cmd{push} するたびに、変更が自動でPRに追加されます}。レビューの指摘を直したら、同じブランチで \cmd{add} → \cmd{commit} → \cmd{push} すればOKです。PRを作り直す必要はありません。

% ---------------------------------------------------------------------
\section{レビューとマージ}

PRの画面では、\textbf{Files changed} タブで変更の差分（追加＝緑、削除＝赤）を確認できます。チーム開発では、ここで他のメンバーがコメントや承認（Approve）を行います。

差分を見て「この変更で問題ないか」を確認します。修正が必要なら作業ブランチに戻って直し、\cmd{push} します。レビューが済んだら、PR画面の \textbf{「Merge pull request」} ボタンで \cmd{main} に取り込みます。

\screenshot[0.62]{03_branch・PR/PR_Change_Program.png}{変更内容の確認とチェック}
\screenshot[0.62]{03_branch・PR/Review_Merge.png}{レビューとマージ}

マージが済んだら、後片付けをします。

\begin{codebox}[bash]
\begin{code}
# 1. ローカルの main に戻る
git switch main

# 2. マージされた最新の main を取り込む
git pull

# 3. 役目を終えた作業ブランチを削除する（ローカル）
git branch -d feature/add-login

# 4. リモートの作業ブランチも削除する（GitHubのPR画面のボタンでも可）
git push origin --delete feature/add-login
\end{code}
\end{codebox}

マージ後にブランチを消すのは、散らかり防止のためです。GitHubのPR画面にも、マージ後に \textbf{「Delete branch」} ボタンが出ます。

次のタスクを始めるときは、また \cmd{main} を最新にしてから新しいブランチを切ります。この\textbf{「最新の \cmd{main} から枝を出す」}を守るのが、コンフリクトを減らす一番のコツです。

% ---------------------------------------------------------------------
\section{実践フロー：ブランチ運用の一周}

ここまでを1つの流れにまとめます。\textbf{この一周を繰り返す}のが、チーム開発の基本リズムです（図\ref{fig:branch-cycle}）。コンフリクトが起きたときの直し方は次章で扱います。

\begin{numfig}{ブランチ運用の一周}{fig:branch-cycle}
\begin{tikzpicture}[node distance=0.75cm]
  \node[flownode, fill=gitGreen, text width=7cm] (A) {git switch main → git pull（最新化）};
  \node[flownode, fill=gitOrange, text width=7cm, below=of A] (B) {git switch -c feature/xxx（枝を出す）};
  \node[flownode, fill=gitRed, text width=7cm, below=of B] (C) {編集 → git add → git commit};
  \node[flownode, fill=gitBlue, text width=7cm, below=of C] (D) {git push -u origin feature/xxx};
  \node[flownode, fill=gitPurple, text width=7cm, below=of D] (E) {GitHubでPR作成};
  \node[flownode, fill=gitBlue, text width=7cm, below=of E] (F) {レビュー → マージ};
  \node[flownode, fill=gitGreen, text width=7cm, below=of F] (G) {main を pull、済んだブランチを削除};
  \draw[flowarrow] (A) -- (B);
  \draw[flowarrow] (B) -- (C);
  \draw[flowarrow] (C) -- (D);
  \draw[flowarrow] (D) -- (E);
  \draw[flowarrow] (E) -- (F);
  \draw[flowarrow] (F) -- (G);
  \draw[flowarrow] (G.east) to[out=0, in=0]
        node[edgelabel, midway, right, xshift=10pt] {繰り返す} (A.east);
\end{tikzpicture}
\end{numfig}

\begin{codebox}[bash]
\begin{code}
# 1. main を最新にする
git switch main
git pull

# 2. 作業ブランチを切る
git switch -c feature/add-login

# 3. 編集して記録する（本編の黄金フロー）
git add .
git commit -m "ログイン機能を追加"

# 4. リモートへ上げる
git push -u origin feature/add-login

# 5. GitHubでPRを作成 → レビュー → マージ（ブラウザで操作）

# 6. 後片付け
git switch main
git pull
git branch -d feature/add-login
\end{code}
\end{codebox}

% ---------------------------------------------------------------------
\section*{この章のまとめ}
\addcontentsline{toc}{section}{この章のまとめ}

\begin{itemize}[leftmargin=1.4em]
  \item \textbf{\cmd{main} は直接触らず}、タスクごとに\textbf{作業ブランチ}を切って編集する
  \item 作業ブランチは \cmd{git push -u origin ブランチ名} でリモートにも共有する
  \item \textbf{PR}は「\cmd{main} に取り込んでほしい」という依頼であり、\textbf{レビューの入口}
  \item マージ後は \textbf{\cmd{main} を \cmd{pull} → 済んだブランチを削除}して片付ける
  \item \textbf{最新の \cmd{main} から枝を出す}ことがコンフリクト予防の要
\end{itemize}

次章では、同じ箇所を別々に変更したときに起きる\textbf{コンフリクト（競合）}の直し方を扱います。

% =====================================================================
% 別冊 第2章 コンフリクト解決
% 元原稿: 03_ブランチ・PR編.md（後半：コンフリクト解決）
% =====================================================================
\chapter{コンフリクト（競合）を解決する}

\textbf{コンフリクト（競合）}は、\textbf{同じファイルの同じ箇所を、2つの流れが別々に変更した}ときに起こります。Gitが「どちらを採用すべきか自動で判断できない」ため、あなたに解決を求めてきます。前章のブランチ運用の中で、マージのときに表面化します。

% ---------------------------------------------------------------------
\section{なぜ起きるのか}

共有リポジトリを元にAさんとBさんが別々に作業し、同じ行を違う内容に書き換えたとします。片方が先にマージした後、もう片方が \cmd{pull} や \cmd{merge} をすると、履歴が食い違ってコンフリクトになります（図\ref{fig:conflict-why}）。

\begin{numfig}{同じ行を別々に変更した2つの流れを合流させるとコンフリクトになる}{fig:conflict-why}
\begin{tikzpicture}
  \node[flownode, fill=gitBlue, text width=3cm] (S) {共有の出発点\\著者: 山田};
  \node[flownode, fill=gitPurple, text width=3.2cm, below left=1.3cm and 0.3cm of S] (A)
        {Aさんが変更\\著者: 山田, 佐藤, 鈴木};
  \node[flownode, fill=gitOrange, text width=3.2cm, below right=1.3cm and 0.3cm of S] (B)
        {Bさんが変更\\著者: 山田, 田中};
  \node[decisionnode, fill=gitAmber, text width=3cm, below=3.4cm of S] (M)
        {同じ行が2通りに…\\どちらが正しい？};
  \node[flownode, fill=gitRed, below=1.2cm of M] (C) {コンフリクト発生};
  \draw[flowarrow] (S) -- (A);
  \draw[flowarrow] (S) -- (B);
  \draw[flowarrow] (A) -- (M);
  \draw[flowarrow] (B) -- (M);
  \draw[flowarrow] (M) -- (C);
\end{tikzpicture}
\end{numfig}

\begin{hintbox}
\textbf{予防が最善です。} 「1タスク＝1ブランチ」「作業前に \cmd{main} を \cmd{pull}」「触るファイルをタスクごとに分ける」を守ると、コンフリクトは大きく減ります。ただしLaTeXの原稿など、同じファイルを複数人で触る場面では起きやすいので、直し方を知っておきましょう。
\end{hintbox}

% ---------------------------------------------------------------------
\section{発生したときの見え方}

\cmd{git pull} や \cmd{git merge} の際に、次のようなメッセージが出ます。

\begin{codebox}[text]
\begin{code}
Auto-merging report.tex
CONFLICT (content): Merge conflict in report.tex
Automatic merge failed; fix conflicts and then commit the result.
\end{code}
\end{codebox}

\cmd{git status} を打つと、\cmd{Unmerged paths} として競合したファイルが表示されます。

\begin{codebox}[text]
\begin{code}
hoge@n171h147 git-practice % git status
On branch main
Your branch is up to date with 'origin/main'.

You have unmerged paths.
  (fix conflicts and run "git commit")
  (use "git merge --abort" to abort the merge)

Unmerged paths:
  (use "git add <file>..." to mark resolution)
        both added:      hoge1.py

no changes added to commit (use "git add" and/or "git commit -a")
\end{code}
\end{codebox}

% ---------------------------------------------------------------------
\section{マーカーの読み方}

競合したファイルを開くと、ぶつかった箇所が次のようなマーカーで囲まれています。ここでは \cmd{report.tex} の著者行を、\textbf{自分は「山田・佐藤・鈴木」}、\textbf{相手は「山田・田中」}に変えてぶつかった例で見ます。

\noindent\textbf{① 発生時の実際の見え方}

\begin{codebox}[text]
\begin{code}
<<<<<<< HEAD
\author{山田, 佐藤, 鈴木}
=======
\author{山田, 田中}
>>>>>>> origin/main
\end{code}
\end{codebox}

\begin{itemize}[leftmargin=1.4em]
  \item \cmd{<<<<<<< HEAD} 〜 \cmd{=======} … \textbf{自分側}（今いるブランチ）の内容 → \cmd{\textbackslash author\{山田, 佐藤, 鈴木\}}
  \item \cmd{=======} 〜 \cmd{>>>>>>> origin/main} … \textbf{相手側}（取り込もうとした内容）→ \cmd{\textbackslash author\{山田, 田中\}}
\end{itemize}

\cmd{HEAD} は「今の自分」を指す目印だと覚えておくと、上が自分・下が相手だと迷いません。

\noindent\textbf{② 手で編集する}

\textbf{どちらを残すか（または両方を活かすか）を決め}、\cmd{<<<<<<<} / \cmd{=======} / \cmd{>>>>>>>} の\textbf{3行をすべて削除}して、正しい1行の状態にします。

\noindent\textbf{③ 解決後の完成形}

たとえば「両方を活かす」なら、最終的にファイルはこうなります。

\begin{codebox}[text]
\begin{code}
\author{山田, 佐藤, 鈴木, 田中}
\end{code}
\end{codebox}

マーカーは1文字も残さず、\textbf{そのまま原稿として通る状態}にするのがゴールです。

% ---------------------------------------------------------------------
\section{解決の手順}

全体の流れは次の通りです。\textbf{発生 → 中身を見て「どれを残すか」決める → マーカーを消して正しい状態に → \cmd{add} → \cmd{commit}}。決めるところで、図\ref{fig:conflict-resolve}のように3択に枝分かれします。

\begin{numfig}{コンフリクト解決の流れ（3択に枝分かれ）}{fig:conflict-resolve}
\begin{tikzpicture}[node distance=0.5cm and 1.4cm]
  \node[flownode, fill=gitRed, text width=2cm] (A) {コンフリクト発生};
  \node[decisionnode, fill=gitAmber, text width=2.2cm, right=of A] (C) {どれを残すか決める};
  \node[flownode, fill=gitPurple, text width=2.6cm, right=of C] (C2) {相手側\\山田, 田中};
  \node[flownode, fill=gitBlue, text width=2.6cm, above=of C2] (C1) {自分側\\山田, 佐藤, 鈴木};
  \node[flownode, fill=gitTeal, text width=2.6cm, below=of C2] (C3) {両方\\山田, 佐藤, 鈴木, 田中};
  \node[flownode, fill=gitOrange, text width=2.2cm, right=of C2] (D) {マーカー3行を削除};
  \node[flownode, fill=gitGreen, right=of D] (E) {git add → commit};
  \draw[flowarrow] (A) -- (C);
  \draw[flowarrow] (C) -- (C1);
  \draw[flowarrow] (C) -- (C2);
  \draw[flowarrow] (C) -- (C3);
  \draw[flowarrow] (C1) -- (D);
  \draw[flowarrow] (C2) -- (D);
  \draw[flowarrow] (C3) -- (D);
  \draw[flowarrow] (D) -- (E);
\end{tikzpicture}
\end{numfig}

解決のやり方は3通りあります。チーム開発ではコンフリクトはPR上で表面化するので、\textbf{PR画面のまま直せる方法A（GitHub上）を基本}にしましょう。手元でじっくり直したいときや、ブラウザで扱えない複雑な競合のときは方法B（VS Code）を、エディタを使わず自分で直したいときは方法C（コマンドライン）を使います。

\subsection{方法A：GitHub（ブラウザ）で解決する（推奨）}

PRを作った後にコンフリクトが起きると、PR画面に \textbf{「This branch has conflicts that must be resolved」} という警告と \textbf{「Resolve conflicts」} ボタンが表示されます。ローカルに戻らず、ブラウザだけで直せる手軽な方法です。

\begin{enumerate}[leftmargin=1.6em]
  \item PR画面の \textbf{「Resolve conflicts」} ボタンを押す。競合したファイルがブラウザ上のエディタで開きます。
\end{enumerate}
\screenshottrim[1.0]{03_branch・PR/conflicts_first_window.png}{コンフリクト発生時のPR画面}{0 240bp 0 0}

\begin{enumerate}[leftmargin=1.6em, start=2]
  \item \textbf{\cmd{<<<<<<<} / \cmd{=======} / \cmd{>>>>>>>} のマーカー}が表示されるので、\textbf{正しい状態に直し、マーカー3行をすべて削除}します。
\end{enumerate}
\screenshottrim[1.0]{03_branch・PR/conflict_fix_file.png}{コンフリクト解決の編集画面}{0 520bp 700bp 0}

\begin{enumerate}[leftmargin=1.6em, start=3]
  \item ファイル上部の \textbf{「Mark as resolved」} を押す。競合ファイルが複数あれば、すべて同様に処理します。
\end{enumerate}
\screenshottrim[1.0]{03_branch・PR/resolve_conflict.png}{コンフリクト解決後のマーク}{0 320bp 0 0}

\begin{enumerate}[leftmargin=1.6em, start=4]
  \item すべて解決したら \textbf{「Commit merge」} を押して確定します。これで競合が解消され、PRをマージできる状態になります。
\end{enumerate}
\screenshottrim[1.0]{03_branch・PR/commit_merge.png}{コンフリクト解決後のコミット}{0 430bp 0 0}

\begin{warnbox}
\textbf{「Resolve conflicts」ボタンが出ないことがあります。} 競合が複雑すぎる場合、GitHubはブラウザでの解決を許可しません。そのときは方法B（VS Code）で、ローカルに取り込んでから解決します。
\end{warnbox}

\subsection{方法B：VS Codeで解決する}

VS Codeは競合を検知すると、専用のボタンで解決を手伝ってくれます。\textbf{マーカーを自分で消す必要はありません。}

\begin{enumerate}[leftmargin=1.6em]
  \item \textbf{競合したファイルを開く。} サイドバーの \textbf{ソース管理（Source Control）} パネルに、競合ファイルが \textbf{「Merge Changes」} としてまとめて表示されます。エクスプローラー上では、そのファイル名の後ろに \cmd{!}（Conflict）マークが付きます。
\end{enumerate}
\screenshottrim[1.0]{03_branch・PR/vscode_conflict_file.png}{VS Codeでのコンフリクト画面（ファイル名の後ろに \texttt{!} が付く）}{0 359bp 0 0}

\begin{enumerate}[leftmargin=1.6em, start=2]
  \item \textbf{ファイルを開くと、「マージエディターで解決」ボタンが表示されます。} この状態で \cmd{<<<<<<<} / \cmd{=======} / \cmd{>>>>>>>} のマーカーを削除し、正しい内容に書き換えることもできます。
\end{enumerate}
\screenshottrim[1.0]{03_branch・PR/vscode_open_conflict_file.png}{VS Codeでのマージエディター起動}{0 359bp 0 0}

\begin{enumerate}[leftmargin=1.6em, start=3]
  \item VS Codeのマージエディターを使うと、\textbf{左＝自分側／右＝相手側／下＝結果}の3ペインで、どちらを採用するかをチェックボックスで選べます。両方を活かして手直ししたいときは、下段を直接編集します。完了したら左側のソース管理パネルで \cmd{+}（ステージ）を押してコミットすれば解決完了です（コマンドの \cmd{git add} → \cmd{git commit} と同じ操作です）。
\end{enumerate}
\screenshottrim[1.0]{03_branch・PR/vscode_merge_editor.png}{VS Codeのマージエディター（3ペイン）}{0 569bp 0 0}

\subsection{方法C：コマンドラインだけで解決する}

エディタを使わずに解決することもできます。ファイルを直接開いてマーカーを消し、\cmd{git add <ファイル>} → \cmd{git commit} で確定します。方法A・Bのボタン操作も、内部ではこれと同じことをしています。

\begin{codebox}[bash]
\begin{code}
# 解決したファイルを「解決済み」として登録する
git add report.tex

# コミットして競合の解決を確定する（メッセージは自動で用意される）
git commit
\end{code}
\end{codebox}

どの方法で解決する場合も、複数ファイルで競合が起きたときは\textbf{すべてのファイルを解決してから}マージを確定します。

\begin{warnbox}
途中で「やっぱりやめたい」ときは \cmd{git merge -{}-abort} で、マージ開始前の状態に戻せます。慌てて壊す前に、まず落ち着いてこのコマンドを思い出してください。
\end{warnbox}

% ---------------------------------------------------------------------
\section*{この章のまとめ}
\addcontentsline{toc}{section}{この章のまとめ}

\begin{itemize}[leftmargin=1.4em]
  \item \textbf{コンフリクト}は同じ箇所の二重変更で起きる
  \item 直し方は \textbf{マーカー3行を消して正しい状態にし、\cmd{add} → \cmd{commit}}
  \item チーム開発では\textbf{GitHubのブラウザ上（方法A）}が基本。複雑なときはVS Code（方法B）
  \item 泥沼化したら \textbf{\cmd{git merge -{}-abort}} で開始前に戻れる
  \item 何より、\textbf{最新の \cmd{main} から枝を出す}ことが予防の要
\end{itemize}

これで、ブランチ運用（第1章）とコンフリクト解決（第2章）が揃いました。チームの運用ルールやPRの書き方は付録Aに、コマンドを忘れたら『逆引きチートシート』、エラーやトラブルに出くわしたら『トラブルシューティング』を辞書代わりに使ってください。


% --- 付録 -------------------------------------------------------------
% トラブルシューティング・チートシートは独立冊子
%（main_trouble.tex / main_cheat.tex）に分離した。
\appendix
% =====================================================================
% 別冊 付録A おすすめのブランチ運用ルールとPRの書き方
% 本文（第1章）で軽く触れた内容の、実践的なテンプレート集。
% =====================================================================
\chapter{おすすめのブランチ運用ルールとPRの書き方}

第1章では、ブランチの切り方やPRの作り方を扱いました。この付録では、\textbf{チームで気持ちよく開発するための「型」}を、そのまま真似できる形でまとめます。絶対の決まりではありませんが、\textbf{迷ったらこれに沿えばまず問題ありません}。チームで最初に決めておくと、後の混乱を防げます。

% ---------------------------------------------------------------------
\section{おすすめのブランチ運用ルール}

\begin{center}
\begin{tabularx}{\linewidth}{@{}>{\raggedright\arraybackslash}p{6.6cm} X@{}}
\toprule
ルール & なぜそうするか \\
\midrule
\cmd{main} は直接編集しない & \cmd{main} を常に「動く・完成した状態」に保つため。GitHubの設定で \cmd{main} への直接 \cmd{push} を禁止することもできる \\
\addlinespace[3pt]
1タスク＝1ブランチ & 目的が1つに絞られ、レビューとマージがきれいにまとまる \\
\addlinespace[3pt]
作業前に最新の \cmd{main} から枝を出す & 古い \cmd{main} から分岐すると、後でコンフリクトが起きやすい \\
\addlinespace[3pt]
こまめにコミット・こまめに \cmd{push} & 事故ってもすぐ戻せる。作業内容を他のメンバーが確認できる \\
\addlinespace[3pt]
PRは小さく保つ & 変更が小さいほどレビューが速く、正確になる \\
\addlinespace[3pt]
マージが済んだらブランチを削除する & 使わないブランチが増えて散らかるのを防ぐ \\
\bottomrule
\end{tabularx}
\end{center}

\medskip

\subsection{ブランチ命名のおすすめ}

ブランチ名は \textbf{\cmd{種類/内容}} の形（英小文字・単語はハイフン区切り）にすると、一覧で目的がひと目で分かります。

\begin{center}
\begin{tabularx}{\linewidth}{@{}>{\raggedright\arraybackslash}p{2.6cm} >{\raggedright\arraybackslash}p{5.4cm} X@{}}
\toprule
接頭辞 & 用途 & 例 \\
\midrule
\cmd{feature/} & 新機能の追加 & \cmd{feature/add-login} \\
\cmd{fix/} & バグ修正 & \cmd{fix/typo-in-readme} \\
\cmd{docs/} & ドキュメントの変更 & \cmd{docs/update-setup} \\
\cmd{refactor/} & 動作を変えない整理 & \cmd{refactor/tidy-utils} \\
\cmd{chore/} & 設定・雑務 & \cmd{chore/update-gitignore} \\
\bottomrule
\end{tabularx}
\end{center}

\medskip

% ---------------------------------------------------------------------
\section{伝わるPRの書き方}

PRは、レビュアーに\textbf{「何を・なぜ変えたか」}を伝える場です。差分（コード）を見れば「どう変えたか」は分かりますが、\textbf{「なぜ」}は書かないと伝わりません。次の点を意識すると、レビューが速く進みます。

\begin{itemize}[leftmargin=1.4em]
  \item \textbf{タイトルは一目で分かるように}。おすすめは \cmd{feat: ログイン機能を追加} のように\textbf{種類を頭に付ける形}（\cmd{feat} 新機能／\cmd{fix} 修正／\cmd{docs} 文書 …）
  \item \textbf{説明には「なぜ」と「背景」を書く}。差分だけでは分からない意図を補う
  \item \textbf{1つのPRは1つの目的に絞る}。ついでの修正を混ぜない
  \item \textbf{動作確認したことを書く}。何をどう確かめたかがあると、レビュアーは安心できる
  \item \textbf{まだ作業中なら Draft（下書き）PR にする}。完成したら「Ready for review」に切り替える
\end{itemize}

% ---------------------------------------------------------------------
\section{PRテンプレート}

毎回ゼロから説明を書くのは大変です。\textbf{PRテンプレート}を用意しておくと、PR作成時に説明欄へ\textbf{自動で見出しが挿入}され、埋めるだけで済みます。

リポジトリの一番上の階層に \cmd{.github/PULL\_REQUEST\_TEMPLATE.md} というファイルを置くと、GitHubがPR作成時に自動で読み込みます。次はそのまま使えるテンプレート例です。

\begin{codebox}[text]
\begin{code}
## 概要
<!-- このPRで何をしたかを1〜2行で -->

## 変更内容
-

## 変更の理由・背景
<!-- なぜこの変更が必要か。差分だけでは分からない意図を書く -->

## 動作確認
- [ ] ビルド / テストが通ることを確認した
- <!-- 確認した手順・結果を書く -->

## 関連 Issue
<!-- あれば #番号。なければ「なし」 -->

## レビューで見てほしい点
<!-- 特に確認してほしい箇所があれば -->
\end{code}
\end{codebox}

テンプレート内の記法について補足します。\cmd{- [ ]} はGitHub上で\textbf{チェックボックス}として表示され、クリックでチェックできます。\cmd{<!-- ... -->} はコメントで、\textbf{記入後は消して}構いません（残しても表示はされません）。

% ---------------------------------------------------------------------
\section{記入例}

上のテンプレートを、\textbf{ログイン機能を追加するPR}で埋めた例です。

\noindent\textbf{タイトル}

\begin{codebox}[text]
\begin{code}
feat: ログイン機能を追加
\end{code}
\end{codebox}

\noindent\textbf{説明}

\begin{codebox}[text]
\begin{code}
## 概要
メールアドレスとパスワードでログインできる画面を追加した。

## 変更内容
- ログイン画面（login.html）を追加
- 入力チェックと認証処理を実装
- ログイン後にトップページへ遷移する処理を追加

## 変更の理由・背景
会員限定ページを用意するため、まず入口となるログイン機能が必要だった。

## 動作確認
- [x] ビルド / テストが通ることを確認した
- 正しいメール・パスワードでログインできることを確認
- 誤ったパスワードでエラーメッセージが出ることを確認

## 関連 Issue
#12

## レビューで見てほしい点
パスワードの扱い方（平文で残していないか）を重点的に見てほしい。
\end{code}
\end{codebox}

\begin{hintbox}
最初から完璧を目指す必要はありません。\textbf{「概要」と「変更の理由」だけでも書いてあると、レビュアーの理解がぐっと速くなります}。チームで運用しながら、自分たちに合う項目に育てていきましょう。
\end{hintbox}


\end{document}
